% SPDX-License-Identifier: Apache-2.0
\section{\code{txhdl\_parts}}
\label{sec:r-parts}

What the language provides as hardware, written in its own lowerable
subset and checked as every unit is. A unit's channel port is one
side of a channel: the lowering makes of it a valid, a ready and the
data, and the testbench that checks the unit plays the channel's
other side from the trace. The channel itself, the elastic buffer of
two between two units, is not in either unit, so two lowered units
wired port to port would have no buffer between them, and their
handshakes would form a combinational loop, which a placer refuses.
\code{Buffer<W>} is the channel as hardware, one per channel of a run
wherever two lowered units meet, on a board or in a larger netlist.

Its rules are the runtime's, from Section~\ref{sec:r-comp}, with no
freedom taken. A take moves the tail into the head. A push goes into
the head if it is empty after that, else into the tail. \code{ready}
on the sending side is the tail being empty as the edge left it,
\code{valid} on the receiving side the head being full, and the data
is the head; all three are registers, so two units on a channel have
no combinational path between them. That the part is the runtime's
channel bit for bit is checked by the examples
document~\cite{examples}, which drives one with a pattern of offers
and takes, does the same on a runtime channel, and asserts after
every cycle that the two agree; the part's netlist is then checked
against that trace under both simulators.

\lstinputlisting[caption={\code{lib/parts/src/buffer.rs}.},label={lst:r-buffer}]{lib/parts/src/buffer.rs}

\code{Fifo<T, AW, D>} is a channel in, a channel out and \code{D}
words between, \code{D = 1 << AW} stated since a computed width in a
type needs nightly Rust. A word is taken from the input whenever
there is room and the oldest word is offered on the output whenever
there is one, so the latency is one cycle; the head and the tail are
\code{AW} bits and a full bit tells full from empty, since the
pointers are equal in both. Behind a unit whose output is a channel
it is the depth the channel's buffer of two does not have. Its test
runs it against a queue.

\lstinputlisting[caption={\code{lib/parts/src/fifo.rs}.},label={lst:r-fifo}]{lib/parts/src/fifo.rs}

\code{Station2} to \code{Station10} are reservation stations, one
type per count of inputs, each written out by \code{station!}
(Section~\ref{sec:r-macros}) since the lowering reads a body and not
a loop over inputs. An input supplies a \code{Tagged<TB, T>}, containing the tag
above the value; the station has \code{L = 1 << TB} lines, a memory
per input for the cells and an occupancy mask per input, one bit per
line. The step is the rule: an input's cell is free when its bit of
the mask is clear; taking an input completes its line when every
other input's cell of that line holds or that input offers the same
tag now; the line sent is the completing input's of lowest index,
when the output has room; an input is taken when its cell is free
and it completes nothing, or completes the line being sent, so an
input never waits on another's cell, only on its own or on the
output's room; a taken input writes its cell and sets its bit, the
line sent clears its bit in every mask, and the values sent come
from the cells or, for the input completing the line, straight from
the input. The tag has two bits at least, since a one-bit value is a
\code{std\_logic} in VHDL and cannot index. The tests run
\code{Station2} against the rule written with loops through four
hundred random cycles, and \code{Station3} on the cycle where two
lines complete at once.

\lstinputlisting[caption={\code{lib/parts/src/station.rs}.},label={lst:r-station}]{lib/parts/src/station.rs}
